% ===============================
% Worksheet / Manual Template
% ===============================
\documentclass[12pt, a4paper]{article}

% ===============================
% Metadata
% ===============================
\newcommand*{\CourseCode}{MAT215}
\newcommand*{\CourseTitle}{Complex Variables \& Laplace Transform}
\newcommand*{\Chapter}{Laplace Transform Using Definition}
\newcommand*{\Topics}{Proof of Formulae and Piecewise Functions}
\newcommand*{\DocDate}{April 08, 2025}

\include{header.tex}

% ==========================================
% Document
% ==========================================
\begin{document}
\FrontPage
\Large

\begin{heading}
	Motivation of Laplace Transform
\end{heading}

\clearpage

\begin{heading}
	Formal Definition of Laplace Transform
\end{heading}

\vspace{10pt}

\begin{definition}[Laplace Transform]
	The \textbf{Laplace transform} of a function $f(t)$, denoted by $\mathcal{L}\{f(t)\}$, is defined as
	\[
		\mathcal{L}\{f(t)\} = F(s) = \int_0^\infty e^{-st} f(t) \, dt,
	\]
	where $s$ is a complex number and the integral converges.
\end{definition}

\clearpage	

\begin{heading}
	Laplace Transform of Common Functions
\end{heading}

\vspace{10pt}

\begin{definition}[Laplace Transform of Algebraic Functions]
	\begin{itemize}
		\item $\mathcal{L}\{1\} = \dfrac{1}{s}$, for $\text{Re}(s) > 0$.
		\item $\mathcal{L}\{t\} = \dfrac{1}{s^2}$, for $\text{Re}(s) > 0$.
		\item $\mathcal{L}\{t^2\} = \dfrac{2}{s^3}$, for $\text{Re}(s) > 0$.
		\item $\mathcal{L}\{t^n\} = \dfrac{n!}{s^{n+1}}$, for $n \in \mathbb{N}$ and $\text{Re}(s) > 0$.
		\item $\mathcal{L}\{t^\alpha\} = \dfrac{\Gamma(\alpha + 1)}{s^{\alpha + 1}}$, for $\alpha > -1$ and $\text{Re}(s) > 0$.
	\end{itemize}
\end{definition}

\vspace{20pt}

\begin{problem}
	Prove that $\mathcal{L}\{1\} = \dfrac{1}{s}$, for $\text{Re}(s) > 0$, using the definition of Laplace transform. 
\end{problem}

\clearpage

\begin{problem}
	Prove that $\mathcal{L}\{t\} = \dfrac{1}{s^2}$, for $\text{Re}(s) > 0$, using the definition of Laplace transform. 
\end{problem}

\clearpage

\begin{problem}
	Prove that $\mathcal{L}\{t^2\} = \dfrac{2}{s^3}$, for $\text{Re}(s) > 0$, using the definition of Laplace transform. 
\end{problem}

\clearpage


\begin{definition}[Laplace Transform of Exponential Functions]
	$\mathcal{L}\{e^{at}\} = \dfrac{1}{s - a}$, for $\text{Re}(s) > \text{Re}(a)$.
\end{definition}

\vspace{10pt}

\begin{problem}
	Prove that $\mathcal{L}\{e^{at}\} = \dfrac{1}{s - a}$, for $\text{Re}(s) > \text{Re}(a)$, using the definition of Laplace transform. 
\end{problem}

\clearpage

\clearpage

\begin{definition}[Laplace Transform of Trigonometric Functions]
	\begin{itemize}
		\item $\mathcal{L}\{\sin(at)\} = \dfrac{a}{s^2 + a^2}$, for $\text{Re}(s) > 0$.
		\item $\mathcal{L}\{\cos(at)\} = \dfrac{s}{s^2 + a^2}$, for $\text{Re}(s) > 0$.
	\end{itemize}
\end{definition}

\vspace{20pt}

\begin{problem}
	Prove that $\mathcal{L}\{\sin(at)\} = \dfrac{a}{s^2 + a^2}$, for $\text{Re}(s) > 0$, using the definition of Laplace transform.
\end{problem}

\clearpage

\begin{problem}
	Prove that $\mathcal{L}\{\cos(at)\} = \dfrac{s}{s^2 + a^2}$, for $\text{Re}(s) > 0$, using the definition of Laplace transform.
\end{problem}

\clearpage

\begin{definition}[Laplace Transform of Hyperbolic Functions]
	\begin{itemize}
		\item $\mathcal{L}\{\sinh(at)\} = \dfrac{a}{s^2 - a^2}$, for $\text{Re}(s) > |a|$.
		\item $\mathcal{L}\{\cosh(at)\} = \dfrac{s}{s^2 - a^2}$, for $\text{Re}(s) > |a|$.
	\end{itemize}
\end{definition}

\vspace{20pt}

\begin{problem}
	Prove that $\mathcal{L}\{\sinh(at)\} = \dfrac{a}{s^2 - a^2}$, for $\text{Re}(s) > |a|$, using the definition of Laplace transform.
\end{problem}

\clearpage

\begin{problem}
	Prove that $\mathcal{L}\{\cosh(at)\} = \dfrac{s}{s^2 - a^2}$, for $\text{Re}(s) > |a|$, using the definition of Laplace transform.
\end{problem}

\clearpage

\begin{heading}
	Laplace Transform of Piecewise Functions Using Definition
\end{heading}

\vspace{20pt}

\begin{problem}
	Using the definition, find the Laplace transform of the function
	\[
	f(t)=
	\begin{cases}
	0, & 0\le t<2,\\
	3, & 2\le t<4,\\
	e^t, & t\ge 4.
	\end{cases}
	\]
		
\end{problem}

\clearpage

\begin{problem}
	Using the definition, find the Laplace transform of the function
	\[
	f(t)=
	\begin{cases}
	0, & 0<t<\dfrac{\pi}{2},\\
	2\cos t, & t\ge \dfrac{\pi}{2}.
	\end{cases}
	\]
	
\end{problem}

\clearpage

\begin{problem}
	Using the definition, find the Laplace transform of the function
	\[
	f(t)=
	\begin{cases}
	\cos 3t, & 0<t<\pi,\\
	0, & t\ge \pi.
	\end{cases}
	\]
		
\end{problem}


\clearpage

\begin{problem}
	Using the definition, find the Laplace transform of the function
	\[
	f(t)=
	\begin{cases}
	0, & 0\le t<\pi,\\
	\sin t, & \pi\le t<2\pi,\\
	0, & t\ge 2\pi.
	\end{cases}
	\]
		
\end{problem}

\clearpage

\begin{problem}
	Using the definition, find the Laplace transform of the function
	\[
	f(t)=
	\begin{cases}
	5\sin t, & 0\le t<\pi,\\
	-4\cos(2t), & t\ge \pi.
	\end{cases}
	\]
		
\end{problem}





\EndPage
\end{document}